\documentclass[10pt,a4paper]{amsart}
\usepackage[margin=19mm]{geometry}
\usepackage{amsmath,amssymb,mathtools}
\usepackage[scr=rsfs,cal=euler]{mathalfa}
\usepackage{xfrac}
\usepackage[unicode,hidelinks,bookmarks=false]{hyperref}
\newcommand{\ie}{i.e.\ }
\newcommand{\cf}{cf.\ }
\newcommand{\resp}{respectively\ }
\newcommand{\Subsection}{\subsection}
\providecommand{\nobreakdash}{\nobreak\hbox{-}}
\setlength{\emergencystretch}{2em}
\multlinegap0pt
\usepackage{amssymb}
\usepackage{xcolor}
\newtheorem{lem}{Lemma}
\newcommand{\bbC}{\mathbb{C}}
\newcommand{\frg}{\mathfrak{g}}
\newcommand{\frh}{\mathfrak{h}}
\newcommand{\frr}{\mathfrak{r}}
\newcommand{\frso}{\mathfrak{so}}
\renewcommand{\ge}{\geqslant}
\renewcommand{\le}{\leqslant}
\newcommand{\wdg}{\wedge}
\title{Rigid currents on compact hyperkahler manifolds}
\author{Nessim Sibony, Andrey Soldatenkov, Misha Verbitsky}
\date{Source: arXiv:2303.11362v3 (CC BY 4.0)}
\begin{document}
\maketitle

\noindent\emph{Source and licence.} Rigid currents on compact hyperkahler manifolds. Version arXiv:2303.11362v3 (CC BY 4.0), distributed by arXiv under the Creative Commons Attribution 4.0 International licence, http://creativecommons.org/licenses/by/4.0/, as recorded in the arXiv licence metadata for this exact version and shown on the abstract page https://arxiv.org/abs/2303.11362v3. Full licence notice: \texttt{licenses/arxiv-2303.11362-CC-BY-4.0.txt}.\par\smallskip

\noindent\textit{Editorial scope.} Counted unit: the article's lem lem\_dim2 (source0.tex lines 1744--1755). The article's complete original proof of it is delivered with it (source0.tex lines 1756--1830, one \texttt{proof} environment of 75 lines). No mathematics is written, repaired or re-derived: the statement and the proof are the article's own lines, in the article's own notation, and no retained line is substituted or re-flowed.  Retained uncounted as the source-local context the proof needs: lines 1676--1678; lines 1680--1685.  Nothing was omitted from any retained span, so the package carries no omission record.  No retained line contains a citation, so the wrapper carries no bibliography and no bibliography entry.  No retained line contains a figure environment, an image-inclusion command or drawing code, so no image is delivered and the package declares no asset dependency.  Editorial formatting only, in addition: an AMS \texttt{amsart} preamble that restates the article's own declarations and macros used by the retained lines, namely the environments \texttt{lem} on separate counters, together with the standard packages those lines need.  The retained environments print in this document as Lemma 1 in order of appearance, whereas the article prints the same items with the numbering of its own full document.  The complete native source of the article as distributed in the arXiv e-print for this exact version is bundled unchanged as \texttt{sources/138940/source0.tex}.
\begin{equation}\label{eqn_set2}
L\subset V_\bbC,\quad\!\! \dim(L) = 2,\quad\!\! 0\neq u\in L^\perp,\quad\!\! q|_L \mbox{ non-degenerate,} \quad\!\! q(u) = 0.
\end{equation}

\begin{lem}\label{lem_bracket}
Assume that $\frh_\bbC\subset\frso(V_\bbC)$ is a Lie subalgebra
that contains $u_1\wdg W_1$ and $u_2\wdg W_2$ for some non-zero isotropic vectors $u_1$, $u_2\in V_\bbC$
and subspaces $W_1\subset u_1^\perp$, $W_2\subset u_2^\perp$.
Assume that there exists $v\in W_1\cap W_2$ with $q(v) \neq 0$. Then $u_1\wdg u_2\in \frh_\bbC$.
\end{lem}

\begin{lem}\label{lem_dim2}
Assume that $d\ge 7$ and that $\frh_\bbC\subset\frso(V_\bbC)$ is a Lie subalgebra
that contains $\frr(L_1,u_1)$, $\frr(L_2,u_2)$ for some $(L_1,u_1)$, $(L_2,u_2)$
as in (\ref{eqn_set2}).
Denote $W_1 = (L_1 + \bbC u_1)^\perp$ and $W_2 = (L_2 + \bbC u_2)^\perp$.
Then one of the following possibilities holds.
\begin{enumerate}
\item either $u_1\wdg u_2\in \frh_\bbC$;
%\item or $\exists v\in u_1^\perp$, $v\notin W_1$ such that $u_1\wdg v\in \frg_\bbC$;
\item or $\exists v\in u_2^\perp$, $v\notin W_2$ such that $u_2\wdg v\in \frh_\bbC$
\end{enumerate}
\end{lem}

\begin{proof}
We will have to consider a number of cases. We note that the codimension of
the subspace $W_1\cap W_2\subset V_\bbC$ is at most 6.

\medskip
{\it Case 1: $\mathrm{codim}(W_1\cap W_2) \le 5$.}

Since $L_1$ is two-dimensional, it contains an isotropic vector $\tilde{u}_1$.
Therefore there exists a two-dimensional isotropic subspace $H_1 = \langle\tilde{u}_1,u_1\rangle\subset (L_1+\bbC u_1)$.
Clearly $H_1\perp W_1$.
First consider the case when the intersection $H_1\cap W_1\cap W_2$ is non-zero,
and denote by $x$ some non-zero vector from this intersection.
Then since $W_1 = (L_1+ \bbC u_1)^\perp$, we have $x\in (L_1+ \bbC u_1)\cap (L_1+ \bbC u_1)^\perp = \bbC u_1$.
It follows that $u_1\in W_2$ and since $u_2\wdg W_2 = \frr(L_2,u_2)$,
we have $u_2\wdg u_1 \in \frr(L_2,u_2)\subset \frh_\bbC$.

So it remains to consider the case when $H_1\cap W_1\cap W_2 = 0$.
We claim that in this case $W_1\cap W_2$ contains a non-isotropic vector.
If not, then $(W_1\cap W_2)\oplus H_1$ is an isotropic subspace in $V_\bbC$.
But this subspace has codimension at most 3, which is strictly less than $d/2$,
hence the subspace cannot be isotropic. Let $v\in W_1\cap W_2$ be a non-isotropic vector.
We apply Lemma \ref{lem_bracket} to conclude that $u_1\wdg u_2\in \frh_\bbC$.

\medskip
{\it Case 2: $\mathrm{codim}(W_1\cap W_2) = 6$.}
We will consider separately two subcases.

\medskip
{\it Case 2a: $q(u_1, u_2)\neq 0$.}

If $u_1\in W_2$ then $u_2\wdg u_1\in u_2\wdg W_2 = \frr(L_2,u_2)\subset \frh_\bbC$. So we may assume
$u_1\notin W_2$, hence $u_1\notin W_2\cap u_1^\perp$.
We claim that the subspace $W_2\cap u_1^\perp$ cannot be isotropic.
If it were, then $(W_2\cap u_1^\perp)+\bbC u_1$ would also be isotropic.
But the latter subspace has codimension at most 3, contradicting the assumption $d\ge 7$.
We conclude that it is possible to find $x\in W_2\cap u_1^\perp$ with $q(x)\neq 0$.

As the next step, we find $y\in (W_1\cap u_2^\perp \cap x^\perp) \setminus W_2$.
To show that it is possible, we need to check that $W_1\cap u_2^\perp \cap x^\perp\nsubseteq W_2$.
Indeed, otherwise we would have
\begin{displaymath}
\mathrm{codim}_{V_\bbC}(W_1\cap W_2) \le \mathrm{codim}_{V_\bbC}(W_1) + \mathrm{codim}_{W_1}(W_1\cap u_2^\perp \cap x^\perp)\le 5,
\end{displaymath}
contradicting our standing assumption in the case we are considering.

We have shown that $\frh_\bbC$ contains the elements $u_1\wdg y$ and $u_2\wdg x$.
Observe that the Lie bracket $[u_1\wdg y, u_2\wdg x]$ is a non-zero multiple of $x\wdg y$.
Next, consider the Lie bracket $[u_2\wdg x,x\wdg y]$. It is a non-zero multiple
of $u_2\wdg y$. Setting $v = y$, we conclude the proof in this case, because
$v$ satisfies the conditions from case (2) in the statement of the lemma.

\medskip
{\it Case 2b: $q(u_1,u_2) = 0$.}

If $u_1\in W_2$, then $u_1\wdg u_2 \in \frr(L_2, u_2)$ and we conclude that $u_1\wdg u_2\in \frh_\bbC$.
So we may assume that $u_1\notin W_2$ and hence $W_2$ is a hyperplane in $W_2+ \bbC u_1$.

We note that $W_2\subset u_1^\perp$ implies that $W_1\cap W_2 = W_2\cap L_1^\perp$, hence
$\mathrm{codim}(W_1\cap W_2)\le 5$, contradicting our assumptions. It follows that
we can find $x\in W_2\setminus (W_2\cap u_1^\perp)$.

Next observe that it is possible to find $y\in (W_1\cap u_2^\perp)\setminus (W_2+ \bbC u_1)$.
Indeed, $W_2$ is a hyperplane in $W_2+ \bbC u_1$. If $W_1\cap u_2^\perp\subset W_2+ \bbC u_1$,
then we intersect $W_2$ and $W_1\cap u_2^\perp$ inside $W_2+ \bbC u_1$, and we see that $W_1 \cap W_2\cap u_2^\perp$ is of
codimension at most one in $W_1\cap u_2^\perp$. Since $W_1$ has codimension 3 in $V_\bbC$,
the codimension of $W_1\cap u_2^\perp$ in $V_\bbC$ is at most 4, hence $\mathrm{codim}_{V_\bbC}(W_1\cap W_2)\le 5$,
contradicting our assumptions.

By construction, $x\in W_2$ and $y\in W_1$, so we obtain two elements $u_1\wdg y$ and $u_2\wdg x$ in $\frh_\bbC$.
Finally, we compute the Lie bracket $[u_1\wdg y, u_2\wdg x] = q(u_1,x)y\wdg u_2 -q(y,x)u_1\wdg u_2
= u_2\wdg(q(y,x)u_1 - q(u_1,x)y)$. Let $v = q(y,x)u_1 - q(u_1,x)y$.
By construction, $q(u_1,x)\neq 0$. Since $y\notin W_2+ \bbC u_1$ we see that $v\notin W_2$.
We see that $v$ satisfies the conditions of case (2) in the statement of
the lemma, and this completes the proof.
\end{proof}

\end{document}
